\section{Introduction}
\label{sec:introduction}

This paper is the technical companion to the conceptual account of Vibe Theory \citep{pollard2026vibeconcept}.
The conceptual paper states what the theory says and how to picture it, deliberately light on
mathematics. This one carries the machinery: the method, the geometry, and a walk through the full
set of findings, eighty-three in all, from a finite and reproducible testbed \citep{pollard2026vibetest}. For
each finding we say what was measured, what came out, and where it stands.

The thesis, stated once, is that reality is one experiential substance, the vibe, and that every
other thing is a pattern of vibes and their relations on a growing, hyperbolic, causal mesh,
updated by one local rule. The conceptual paper defends that picture. Here we test it. The method
is to turn each load-bearing question into a runnable measurement, a finite seeded program that
returns a number and that could fail, guarded by a known-answer test suite. What survives is
reported, sorted into what is demonstrated, what is a first rung on a deep problem, and what is
open.

The substrate is the heart of the model, and it is the heart of this paper. More than anything
else, Vibe Theory rests on a particular kind of geometric object, a regular tessellation of a
negatively curved space, a Coxeter honeycomb, grown rather than placed. Section~\ref{sec:geometry}
treats that object in full, because the natural intuitions about it are mostly wrong and because
the rest of the physics depends on getting it right. We cover the Coxeter reflection groups and
their Schl\"afli symbols, the regular tilings of the hyperbolic plane, the regular honeycombs of
hyperbolic space, the specific crystals the model uses, the classification that pins the spatial
dimension, the reason curvature buys Lorentz invariance, and the deterministic growth that makes
the crystal a process rather than a static orbit. Two of the model's crystals are shown directly:
the heptagrid $\{7,3\}$ and the three-dimensional dodecagrid $\{5,3,4\}$.

The paper is organised bottom-up. Section~\ref{sec:testbed} describes the testbed and the
known-answer discipline. Section~\ref{sec:geometry} is the geometry of the substrate.
Section~\ref{sec:base} shows how that geometry is forced from the integers. Section~\ref{sec:dynamics}
puts the tones and the rule on the substrate and recovers time and energy.
Sections~\ref{sec:matter} through \ref{sec:cosmology} walk matter and forces, gravity, the quantum,
and cosmology. Section~\ref{sec:standardmodel} covers the Standard Model results and what they
assume. Section~\ref{sec:interior} covers the measurable parts of the interior, recursion and
integration and selves. Section~\ref{sec:capstone} runs the committed model end to end, and
Section~\ref{sec:scoreboard} is the full scoreboard.

Throughout, the work sits in the family of discrete-physics programs: the causal sets of Bombelli,
Lee, Meyer and Sorkin \citeyearpar{bombelli1987spacetime}, the lattice gauge theory of Wilson \citeyearpar{wilson1974confinement}, the overlap fermions of
Neuberger \citeyearpar{neuberger1998exactly}, the cellular-automaton interpretation of 't Hooft \citeyearpar{thooft2016cellular}, and, for the geometry,
the regular hyperbolic honeycombs of Coxeter \citeyearpar{coxeter1973regular} and the hyperbolic cellular automata of
Margenstern \citeyearpar{margenstern2013small}. What Vibe Theory adds is the insistence that these live on one experiential
substrate, and the demonstration, in a finite simulation, that they can. References to physics are
sources of structure and tests of coherence, not claims to have derived nature.
